\section{Introduction}\label{sec:introduction}

Longitudinal and clustered studies increasingly record a large number of candidate predictors for each observational unit. Examples include repeated transcriptomic measurements, electronic-health-record trajectories, panel data, and longitudinal imaging features. In such studies, the number of measured predictors may be comparable to, or much larger than, the number of independent subjects. At the same time, observations from the same subject share persistent heterogeneity that cannot be represented adequately by a pooled regression. Linear mixed-effects models address within-cluster dependence through low-dimensional cluster-specific effects \cite{pinheiro2000,jiang2007}, but standard likelihood theory is primarily designed for moderate-dimensional mean regression.

Quantile regression provides a complementary description of heterogeneous associations across the response distribution \cite{koenker1978,koenker2005}. Sparse high-dimensional quantile regression has well-developed estimation theory \cite{belloni2011,tan2022}, and convolution smoothing makes the check loss differentiable while retaining the target of quantile regression up to a controlled smoothing error \cite{fernandes2021,he2021}. Recent work has also established debiased inference for smoothed high-dimensional quantile regression with independent observations \cite{yan2023}. These developments do not directly resolve the geometry induced by cluster-specific nuisance effects.

Mixed-effects quantile regression has been studied through asymmetric-Laplace working likelihoods and penalised formulations \cite{geraci2014,geraci2014jss,li2020}. A central difficulty appears in the many-small-cluster regime: the number of clusters tends to infinity while each cluster contains only a bounded number of observations. The cluster effects then form a growing collection of incidental nuisance parameters. Existing estimators can exhibit material fixed-effect bias in this regime, even in low dimension \cite{battagliola2022}. A valid high-dimensional inferential procedure must therefore identify its population target, profile the incidental effects coherently, and use the resulting profile curvature rather than the curvature of a pooled model.

A parallel literature treats high-dimensional correlated outcomes without explicitly profiling subject effects. Penalised estimating equations and quadratic inference functions have been developed for longitudinal quantile regression and variable selection \cite{zu2023,song2024,bhattacharya2026}. High-dimensional testing for longitudinal data is also available in mean and quantile settings \cite{fang2020,wang2023}. For Gaussian mixed models, \citeA{li2022} established estimation and inference by a quasi-likelihood construction. These methods clarify the importance of cluster-level information, but they do not supply the profile score and effective information needed when low-dimensional cluster-specific effects enter a quantile criterion directly.

This paper develops a profile-inference framework for that setting. The formulation is deliberately centred on the inferential geometry rather than on screening or optimisation as separate contributions. Each cluster is assigned a strictly convex ridge-regularised nuisance subproblem. The population target is defined from the unsmoothed profiled check loss and is therefore independent of the smoothing bandwidth. The complete ridge criterion is then profiled before differentiation. This order matters: the profile score is formed with the original fixed-effect design, whereas the Hessian is a Schur complement. When the Hessian is expressed using a residualised design, an additional positive ridge-curvature term is required. Omitting that term makes the score, Hessian, and one-step correction correspond to different objective functions.

The paper makes three main contributions. First, we define an unsmoothed regularised profile parameter for high-dimensional quantile regression with many small clusters. The target remains well defined when individual cluster effects cannot be consistently estimated, and its distinction from a structural conditional-quantile coefficient is explicit. Convolution smoothing is treated as a computational and curvature-estimation device, with a separate approximation result linking the smoothed and unsmoothed targets.

Second, we derive the exact profile score and profile Hessian for the complete ridge-regularised cluster criterion. The Hessian is estimated cluster by cluster and inverted only in the rows needed for inference through a Dantzig/CLIME program. A one-step correction then yields a cluster-level influence representation. The construction retains within-cluster dependence through empirical cluster scores and does not replace the profile score by a residualised-design score.

Third, we formulate a theoretical regime in which penalised estimation, effective-Hessian estimation, row-wise precision estimation, debiasing, and cluster-sandwich variance estimation are jointly valid. The conditions are stated at the level of independent clusters. An explicit bandwidth choice and sparsity regime is provided to demonstrate that the conditions can hold simultaneously. The numerical protocol prohibits truth-based coverage calibration and evaluates both correctly specified and misspecified working random-effect structures.

The empirical study is redesigned around GSE65391, a longitudinal paediatric systemic lupus erythematosus (SLE) study with substantially more subjects and repeated visits than the dataset used in the earlier manuscript. Whole-blood expression measurements serve as high-dimensional predictors, and complement C3 is used as a continuous response. Lower conditional quantiles of C3 are scientifically relevant because they describe visits with the strongest complement depletion. The analysis combines subject-specific random effects, always-included demographic and treatment variables, and sparse transcriptomic effects.

The remainder of the paper is organised as follows. Section~\ref{sec:model} defines the model, profile target, smoothed estimator, and exact profile calculus. Section~\ref{sec:inference} constructs the debiased estimator and cluster-robust variance. Section~\ref{sec:theory} states the theoretical results and a feasible asymptotic regime. Section~\ref{sec:simulation} specifies the simulation study. Section~\ref{sec:application} gives the complete GSE65391 analysis protocol. Section~\ref{sec:discussion} concludes. Proof details, computational diagnostics, and extended numerical designs are assigned to the \supp.

